% LSPV-2 REPORT::GHWQUAD1-R1::2026-10-04 22:24 PT
\documentclass[11pt]{article}
\usepackage[margin=1in]{geometry}
\usepackage{amsmath,amssymb}
\usepackage{hyperref}
\title{LSPV-2 GHWQUAD1-R1: Generalized-weight certificates for TASK-001}
\author{LSPV-2 separator-algebra/product-span}
\date{2026-10-04 22:24 PT}
\begin{document}
\maketitle

\section*{Status}
P versus NP remains OPEN.  This report records a finite-dimensional structural theorem for the TASK-001 quadratic evaluation operator.  It does not provide the arbitrary-SAT-algorithm return bridge.

\section*{CORPUS REVIEWED}
Repository governance in AGENTS.md and CODEX.md, the repository pnp-research skill, root main.tex, TASK-001, the current LSPV-2 lineage including QUOTIENT1, RETURNSQUEEZE1, DELAYOR1, RECURCLOSE1, SELECTCHARGE1, CONSTQUOT1, and AFFQUAD1, plus the current shared handoff containing LSPV-3 CYCLEGAIN1.  The total tree was surveyed before this recovery write.

\section*{DUPLICATION / CONTRADICTION AUDIT}
The result below strictly extends the earlier minimum-distance density certificate conceptually: the minimum-distance case is the first generalized Hamming weight.  It does not replace AFFQUAD1.  The constant-aware normalization correction remains necessary: collapse equal base labels, retain one zero label when present, and only then form the constant-augmented degree-two evaluation operator.  No source-specific SU/vMS asymptotic rank theorem is claimed.

\section*{Setup}
Let
\[
 V_r=\{p:\mathbb F_2^r\to\mathbb F_2:\deg p\le 2\},
 \qquad
 D_r=\dim V_r=1+r+\binom r2 .
\]
For a distinct-label set \(X\subseteq\mathbb F_2^r\), let
\[
 \rho(X)=\operatorname{rank}(\operatorname{ev}_X),\qquad
 \kappa(X)=D_r-\rho(X),
\]
and write \(Y=\mathbb F_2^r\setminus X\).  Full-cube evaluation of \(V_r\) is the binary second-order Reed--Muller code \(C_r=\mathrm{RM}(2,r)\).

\section*{GHWKERNEL1}
\textbf{PROVED/EXACT.}
Let \(d_j(C_r)\) be the \(j\)-th generalized Hamming weight of \(C_r\).  For every \(j\ge 1\),
\[
 |Y|<d_j(C_r)\quad\Longrightarrow\quad
 \kappa(X)<j,
\]
equivalently
\[
 |Y|<d_j(C_r)\quad\Longrightarrow\quad
 \rho(X)\ge D_r-j+1.
\]

\paragraph{Proof.}
The kernel
\[
 K_X=\{p\in V_r:p|_X=0\}
\]
evaluates injectively on the full Boolean cube to a subcode of \(C_r\) supported inside \(Y\).  Its dimension is exactly \(\kappa(X)\).  If \(\kappa(X)\ge j\), choose a \(j\)-dimensional subcode of that image.  Its support is contained in \(Y\), so by the definition of generalized Hamming weight,
\[
 d_j(C_r)\le |Y|,
\]
contradicting the hypothesis. \(\square\)

For \(j=1\), the standard minimum-distance identity
\[
 d_1(\mathrm{RM}(2,r))=2^{r-2}
\]
recovers the earlier density certificate
\[
 |X|>\tfrac34 2^r\Longrightarrow \rho(X)=D_r.
\]

\section*{Affine-local form}
\textbf{PROVED/EXACT.}
Let \(W\subseteq\mathbb F_2^r\) be an affine subspace of dimension \(s\), and let \(Y_W=W\setminus X\).  By affine invariance of degree-at-most-two Boolean evaluation,
\[
 |Y_W|<d_j(\mathrm{RM}(2,s))
 \Longrightarrow
 \rho(X)\ge 1+s+\binom s2-j+1.
\]
Thus a sufficient condition for the TASK-001 structural target is the existence of affine flats \(W_m\) and indices \(j(m)\) such that
\[
 |W_m\setminus X_m|<d_{j(m)}(\mathrm{RM}(2,s(m)))
\]
and
\[
 1+s(m)+\binom{s(m)}2-j(m)=\omega(m).
\]
The implication is PROVED/EXACT; whether the exact source-transcribed SU/vMS family satisfies these premises is OPEN.

\section*{Source-specific TASK-001 boundary}
The exact SU/vMS Boolean point set \(X_m\) is still missing the complete source transcription demanded by TASK-001: exact generator theorem, field/vector-space identification, binary encoder, Boolean projections, local query algorithm, adaptive dependencies, and final output-coordinate definition.  Therefore no schematic Reed--Muller or orbit-window surrogate may be substituted.  The next structural computation should be the affine-flat generalized-weight profile only after those source definitions are complete.

\section*{Return-bridge accounting}
GHWKERNEL1 is a structural rank certificate only.  It does not force an arbitrary polynomial-time SAT decider to acquire \(X\), its quadratic kernel, or a generalized-weight witness.  SELECTCHARGE1 remains active: unresolved calls to the original decider on current-scale residuals retain dependence on the unknown starting exponent.  CYCLEGAIN1 refines the recursive accounting: an individual scale-one typed edge need not be fatal when recursion is well-founded and every complete directed cycle has multiplicative scale gain below one, but a direct unresolved current-scale call to the original decider still carries the original exponent and materialized scale-one queries retain their workspace charge.

\section*{Hostile controls and barrier audit}
COPY does not change the distinct-label set.  DEAD requires retaining one zero label because the constant coordinate is nonzero there.  Invertible affine INTERP preserves the quadratic evaluation problem.  Nonlinear INTERP and RECOMP remain live attacks.  Easy affine/Tseitin, Horn, 2-SAT, bounded-treewidth, and independent-product families can also exhibit large quadratic evaluation rank, so rank alone is not a SAT-hardness certificate.  The theorem is algebraic/coding-theoretic and does not by itself cross relativization or algebrization barriers.

\section*{PQC transfer classification}
Generalized-weight support accounting and Schur-square/shortening diagnostics are EXACT OBJECT/THEOREM TRANSFERS once the code is fixed.  Their use as an unrestricted SAT lower bound is a FAILED TRANSFER absent a source-matched acquisition/return theorem.  Shortening and quadratic-relation spectra remain ANALOGY/CANDIDATE diagnostics for locating structure in the exact source family.

\section*{WHAT CHANGED RELATIVE TO ROOT MAIN}
The structural ledger can be sharpened from an all-or-nothing full-rank test to a graded kernel bound controlled by generalized Hamming weights.  Root normalization language should also remain constant-aware.  No executive P-versus-NP status changes.

\section*{ADOPTED / REJECTED / RETIRED PRIOR COMPONENTS}
Adopt AFFQUAD1 affine invariance and the constant-aware evaluation formulation.  Adopt SELECTCHARGE1 and CYCLEGAIN1 accounting guards.  Retain QMINOR1 as a sufficient certificate but not a necessary one.  Reject any inference from high quadratic rank directly to unrestricted SAT runtime.

\section*{STALE OR OVERSTATED CLAIMS FOUND}
Any instruction to delete zero base columns before forming the constant-augmented TASK-001 operator is stale.  Any statement that every individual scale-one typed edge automatically blocks recursive exponent erasure is too strong after CYCLEGAIN1; the cycle condition and workspace gate must be stated separately.

\section*{Next target}
Complete the exact source transcription required by TASK-001.  Then compute the affine-flat generalized-weight profile of the resulting distinct constant-aware Boolean point set \(X_m\), seeking either an \(\omega(m)\) certified rank lower bound or an explicit low-dimensional quadratic-relation obstruction.

\section*{Final status}
P versus NP remains OPEN.  GHWKERNEL1 and its affine-local corollary are PROVED/EXACT finite-dimensional structural statements.  The source-specific SU/vMS premise and the arbitrary-algorithm return bridge remain OPEN.

% END_OF_VERIFIED_REPORT::LSPV-2::GHWQUAD1-R1
\end{document}
